\scanpage{24}
% @ 1–3, top
donc définit un morphisme \deleted{canonique}
\[
X\longrightarrow X_0
\]
qui est surjectif radiciel fini.

% @ 3–12, middle
Pour montrer que ce morphisme est bien solution du \unsure{Pb} universel que nous avions dit, on voit que sur $Y^{*}$, un morphisme de $X^{*}$ dans un \unsure{revêtt|revêtement} étale $Y'{}^{*}$ de $Y^{*}$ se factorise de façon unique en
\[
X^{*}\unsure{-|\to} X_0^{*}\longrightarrow Y'{}^{*}
\]
[et si $Y'{}^{*}$ provient d'un $Y'$, alors le morphisme donné est compatible avec les données de descente \deleted{\unsure{ssi}} $X^{*}_0\to Y'{}^{*}$ l'est] Pour vérifier ce fait, on est ramené aussitôt au cas où $X^{*}_0=Y^{*}$, où c'est immédiat \unsure{---} (\unsure{vu -- lemme|ou -- lemme|cf. lemme}).

% @ 16–19, lower paragraph
\uline{À \deleted{\unsure{vérifier}}} \uline{généraliser} au cas où on a un morphisme \unsure{propre} \uline{univ\textsuperscript{t} ouvert} $X\to Y$ dans lequel les composantes connexes \inserted{absolues} des fibres ``se comportent bien'' i.e. en nb \inserted{loc} constant \deleted{\unsure{et}}

% @ 20–21, foot
\uline{Question} Caractérisation des foncteurs représentables par des rev. étales ???
% @ middle, watermark and reverse letterhead
% printed form
\textsc{Extra Strong}\\
\unsure{Guerimand|Guérimand} Voiron
% The reversed letterhead at upper right shows through from the other face.
% archive sheet number: 23
